\chapter{Background: Dose-Response Modeling, Metrology, and Activation Steering}\label{ch02:background}

\section{Purpose, scope, and status of this chapter}\label{ch02:purpose}

This chapter is a tutorial survey of the established literatures on which CAS-Ohm 1.0 stands. It covers, in order: dose-response modeling, the empirical tradition from which the candidate takes its analytical frame (Section~\ref{ch02:dose}); metrology, the science of measurement, from which the candidate takes its quantity-and-uncertainty discipline (Section~\ref{ch02:metrology}); the statistical calibration of computer models, from which the candidate takes its treatment of fitted parameters, model discrepancy, and split discipline (Section~\ref{ch02:calibration}); verification and validation of models, from which the candidate takes the distinction between tested software and validated science (Section~\ref{ch02:vv}); the limits of causal inference from observational fits, which bound what any fitted response relation can ever mean (Section~\ref{ch02:causal}); activation steering and representation engineering in neural networks, which constitute the empirical setting in which the candidate's quantities would be measured (Section~\ref{ch02:steering}); and provenance modeling, from which the candidate takes its lineage vocabulary (Section~\ref{ch02:prov}). A closing synthesis (Section~\ref{ch02:synthesis}) states what this background licenses for the rest of the monograph and, just as importantly, what it does not.

Two status rules govern every sentence that follows, and the reader is asked to hold the chapter to them.

\begin{ruleenv}[Background is not novelty]\label{ch02:rule:notnovel}
Nothing in this chapter is a contribution of the CAS-Ohm project. Every concept, method, distinction, and caution surveyed here belongs to its own literature and is presented as such. Where later chapters adopt one of these concepts, the adoption is an application of standard discipline, never a claim of invention.
\end{ruleenv}

\begin{ruleenv}[Background smuggles no empirical claims]\label{ch02:rule:nosmuggle}
Nothing in this chapter asserts, implies, or presupposes any empirical result about CAS-Ohm 1.0. Descriptions of steering experiments, dose-response regions, calibration successes, or provenance systems are descriptions of the external literature and of generic methodology. In particular, the survey of activation steering in Section~\ref{ch02:steering} describes what that literature does and reports at primer depth; it does not state or suggest that any CAS-Ohm quantity has been measured, that any CAS-Ohm hypothesis has support, or that the phenomena the literature reports transfer to the candidate's setting. All CAS-Ohm claims remain \textsc{open / ring 4 / not licensed / held}.
\end{ruleenv}

The chapter is deliberately elementary in places. The candidate's historical sources went wrong not on advanced questions but on elementary ones --- what a unit is, what a fitted parameter can claim, what a residual can prove --- and the corrective is a survey that states the elementary things carefully, with citations to the authoritative treatments, so that later chapters can invoke them by reference rather than re-arguing them.

\section{Dose-response modeling}\label{ch02:dose}

\subsection{The setting}\label{ch02:dose:setting}

Dose-response analysis is the branch of applied statistics concerned with how a measured response varies as a function of a controlled stimulus, conventionally called the dose. The tradition is oldest and most developed in the biological sciences --- pharmacology, toxicology, weed science, environmental risk assessment --- where the dose is a chemical concentration or exposure and the response is growth, mortality, binding, or some other measured endpoint, but its formal apparatus is domain-general: a design variable swept over a grid, a response observed with error, and a parametric or semi-parametric curve fitted to summarize the relation \cite{ritz}.

Three features of the tradition matter for this monograph. First, dose-response analysis is unapologetically \emph{empirical}: the fitted curve is a summary of observed behavior under a protocol, not a derivation from first principles, and its parameters --- slopes, midpoints, asymptotes --- are estimates with standard errors, not constants of nature. Second, the tradition is organized around \emph{characteristic regions} of the dose axis, described next, which give qualitative structure to an otherwise arbitrary curve. Third, the tradition has accumulated hard-won cautions about extrapolation, model selection, and design --- cautions that transfer directly to any new domain that borrows the frame.

\subsection{Model families}\label{ch02:dose:families}

The workhorse models of dose-response analysis are sigmoidal. The log-logistic family, in its four-parameter form, expresses the mean response at dose $x$ as
\[
f(x) = c + \frac{d - c}{1 + \exp\bigl(b\,(\log x - \log e)\bigr)},
\]
with lower asymptote $c$, upper asymptote $d$, steepness parameter $b$, and midpoint (the dose of half-maximal effect) $e$; specializations fix $c = 0$ or $d = 1$, and relatives such as the Weibull and log-normal families trade tail behavior against the same qualitative shape \cite{ritz}. What all of these share is the sigmoid's three-act structure: a low-dose plateau, a rising middle, and a high-dose plateau. Threshold variants add an explicit dose below which the mean response is exactly at baseline; hormetic variants permit non-monotone behavior at low dose. The family chosen constrains what the analyst can express: a model without an upper asymptote cannot represent saturation, and a model without a threshold parameter will interpolate smoothly through what may in truth be a hard onset.

Selection among families is itself a statistical act with consequences. Information criteria, lack-of-fit tests against richer alternatives, and comparison across nested specializations are standard practice, and the literature is insistent that the fitted family is a modeling choice to be reported and defended, not a fact discovered in the data \cite{ritz}. An analyst who fits only a straight line and reports a slope has not shown the relation is linear; an analyst who fits a sigmoid and reports a midpoint has not shown saturation exists. The curve constrains the summary; the data constrain the curve only within the family offered.

\subsection{Response regions}\label{ch02:dose:regions}

The qualitative anatomy of a sigmoidal dose-response relation is conventionally described in regions, and this anatomy --- not any electrical analogy --- is the ancestral home of the region vocabulary used later in this monograph.

\emph{Sub-threshold region.} At sufficiently low dose, the response is statistically indistinguishable from baseline: whatever mechanism couples dose to response has not been engaged above the noise floor. Operationally, the boundary of this region is a detection question as much as a mechanism question --- it depends on assay noise, replication, and the baseline's own variability --- and the literature distinguishes carefully between ``no effect'' and ``no detected effect.''

\emph{Approximately linear region.} In the rising middle of a sigmoid, the curve is locally well-approximated by a straight line, and for many designs the analyst's primary quantity is exactly the local slope in this region: the marginal response per unit dose. The linear region is where proportional summaries are honest, where relative potencies of two treatments can be compared by slope ratios, and where the classical machinery of linear calibration applies. Its boundaries are model-dependent estimates, not observables, and reporting a slope without reporting the region over which linearity was assessed is a recognized malpractice.

\emph{Saturation region.} At high dose the response approaches its upper asymptote: receptors are occupied, transport is maximal, the instrument clips, or the underlying process has no further room to respond. Marginal response per unit dose falls toward zero, and two consequences follow that the literature emphasizes. First, doses in the saturation region are nearly uninformative about the slope parameter --- the design's information about proportionality lives in the middle. Second, summaries computed in saturation are treacherous: a ratio of response to dose declines mechanically with dose even when the underlying relation is perfectly regular, so a declining ratio is not evidence of a changing mechanism.

\emph{Beyond saturation.} Some assays exhibit qualitative breakdown at extreme dose --- toxicity supervening on the measured endpoint, instrument overload, loss of the preparation --- where the response is no longer a measurement of the original quantity at all. The honest treatment is exclusion with a recorded reason, not inclusion as data.

Two further phenomena complete the anatomy. \emph{Non-monotone} relations (hormesis in the biological literature) violate the sigmoid's ordering assumption outright and demand richer families. And \emph{instability} --- irreproducibility of the response at fixed dose across replicates beyond assay noise --- signals that the preparation, not the curve, is the problem; no fitted parameter from an unstable preparation deserves trust.

\subsection{Estimation and inference}\label{ch02:dose:estimation}

Fitting a dose-response model is, computationally, nonlinear regression: the family's parameters are estimated by minimizing a loss --- ordinarily least squares under an assumption of independent errors, or a transformation- or weight-adjusted variant when the response variance changes with its mean, as it commonly does near asymptotes \cite{ritz}. Several inferential practices from this setting are recorded here explicitly, because each has a disciplinary echo later in the monograph.

First, \emph{parameter uncertainty is part of the estimate}. The fitted slope, midpoint, and asymptotes come with standard errors derived from the curvature of the loss surface (or from resampling, when the asymptotic approximations are doubtful --- bootstrap intervals are routine in the applied literature), and downstream quantities computed from the parameters inherit uncertainty through the delta method or through resampling. A reported midpoint without an interval is regarded as an incomplete result, not a bold one.

Second, \emph{variance structure is a modeling decision}. Dose-response data are frequently heteroscedastic --- noisier in the rising region than at the plateaus, or vice versa --- and ignoring this does not merely widen intervals; it can bias which region dominates the fit. Weighted fitting, variance modeling, and transform-both-sides approaches are the standard responses, and the choice made is reported with the model \cite{ritz}.

Third, \emph{replicates and lack-of-fit are the honesty check}. With genuine replicates at each dose, the pure-error variance is estimable independently of any curve, and a lack-of-fit comparison asks whether the chosen family's residual variation exceeds what replicates alone explain. A family that fails lack-of-fit is not rescued by the elegance of its parameters; it is replaced or the failure is reported.

Fourth, \emph{designs are evaluated by the information they carry about the target parameter}. If the analysis's purpose is a slope in the linear region, doses concentrated at the plateaus are wasted; if the purpose is an asymptote, the reverse. Optimal-design considerations --- placing doses where the parameter of interest is most informed --- are part of the tradition's standard toolkit, and their qualitative lesson (the grid is chosen for the question, and a grid chosen for another question may be uninformative for this one) transfers to any borrowed setting.

\subsection{Derived quantities: effective doses and relative potency}\label{ch02:dose:derived}

The tradition's headline summaries are typically \emph{derived} quantities: the effective dose $\mathrm{ED}_{50}$ (the dose producing half-maximal response, coinciding with the midpoint parameter in the standard families), its generalizations $\mathrm{ED}_p$ for other fractions, and the \emph{relative potency} of two treatments, defined --- under an assumption of similar curve shape --- as the ratio of doses producing equal response \cite{ritz}. Two cautions attach. Derived quantities inherit and often amplify parameter uncertainty: an $\mathrm{ED}_p$ far from the sampled region can have an interval spanning orders of magnitude, and reporting the point value without the interval misrepresents the experiment. And comparability constructs like relative potency are \emph{conditional on shape similarity}: when two curves differ in slope or asymptote, a single potency ratio does not exist, and forcing one manufactures a number without a referent. The general lesson --- a derived summary exists only under the assumptions of its derivation, and those assumptions are checkable and reportable --- is applied repeatedly in later chapters, where every derived CAS-Ohm quantity carries its derivation parents and method identity precisely so that the conditioning assumptions remain visible.

\subsection{Design and the flat-response ambiguity}\label{ch02:dose:design}

A methodological staple of the tradition is that the dose grid is part of the inference. A grid confined to low doses cannot distinguish a sub-threshold relation from an insensitive assay; a grid confined to high doses cannot distinguish saturation from a ceiling artifact; a two-point grid cannot detect curvature at all. Good designs bracket the anticipated midpoint, replicate at each dose, and include controls at zero dose.

Of particular relevance to later chapters is what may be called the \emph{flat-response ambiguity}: an observed response that does not vary with dose is compatible with at least four distinct states of affairs --- the dose is sub-threshold; the response is saturated; the assay is insensitive or mis-targeted; or the measurement failed. Distinguishing them requires information beyond the flat curve itself: positive controls, grid extension, assay validation. The dose-response literature treats this as elementary; it is recorded here because the corresponding conflation (reading every flat response as one particular failure) is precisely the mnemonic overload that the source audit found in the candidate's lineage, and the disposition there --- classify the dose-response signature and the semantic outcome independently --- is an application of this standard point, not an innovation.

\subsection{Calibration curves and working ranges}\label{ch02:dose:calibcurve}

A sibling of the dose-response curve, from analytical measurement, is the \emph{calibration curve}: the fitted relation between an instrument's indication and the known values of reference materials, used thereafter to convert indications into reported values. The two curves are formally similar --- a controlled abscissa, a measured ordinate, a fitted relation --- but their inferential directions differ: dose-response analysis \emph{characterizes} the relation as the object of interest, while calibration-curve practice \emph{inverts} it, reading unknown inputs off measured outputs. The inversion imports every property of the fitted curve into every subsequent measurement, which is why analytical practice surrounds calibration curves with explicit range discipline.

Three range concepts from that practice complete the vocabulary this monograph borrows. The \emph{limit of detection} is the smallest input distinguishable from blank with stated confidence --- below it, indications are reported as ``not detected,'' never as small numbers. The \emph{limit of quantification} is the smallest input measurable with stated relative uncertainty --- between detection and quantification, presence may be asserted but magnitude may not. The \emph{working range} (or linear range, when linearity is the fitted form) is the interval over which the calibration relation has been established with acceptable lack-of-fit --- and reported values outside it are extrapolations of the calibration, flagged as such or refused. Together these define a regime discipline: an instrument's conversion relation is trusted only inside a characterized envelope, with distinct sub-thresholds for ``say nothing,'' ``say present,'' and ``say how much.''

The relevance to a response-model candidate is structural. Any proposal that fits a relation between an intervention coordinate and a response coordinate, and then uses the fitted parameter to predict, is running calibration-curve logic; it thereby inherits the obligations of that logic --- established range, declared limits below which no quantitative claim is made, refusal to report outside the envelope --- independently of any question about mechanism. The candidate's cutoff and saturation gating and its refusal to emit diagnostics from out-of-region operating points, described in later chapters, are these obligations transcribed; the obligations themselves are as old as analytical chemistry.

\subsection{What the tradition offers this monograph}\label{ch02:dose:offer}

From dose-response analysis, the candidate inherits: the framing of a controlled scalar dose and measured scalar response; the region anatomy (sub-threshold, approximately linear, saturating, broken-down, non-monotone, unstable) as a classification of \emph{signatures} rather than mechanisms; the insistence that linearity is a locally assessed, design-dependent, reported property rather than an assumption; the caution that ratios computed outside the linear region mislead; and the flat-response ambiguity. What the tradition does not offer --- and what none of this chapter should be read as offering --- is any evidence that steered neural activations exhibit these anatomies. That is an open empirical question, registered as such.

\section{Metrology: quantities, units, uncertainty, traceability}\label{ch02:metrology}

\subsection{Why metrology, and which metrology}\label{ch02:metrology:why}

Metrology is the science of measurement and its applications. Its two foundational reference documents are the International Vocabulary of Metrology, which fixes the concepts and terms --- quantity, unit, measurand, calibration, traceability --- \cite{vim}, and the Guide to the Expression of Uncertainty in Measurement, which fixes the treatment of measurement uncertainty \cite{gum}. The candidate borrows this conceptual discipline wholesale while claiming none of its institutional substance: CAS-Ohm has no accredited laboratories, no reference standards, no interlaboratory comparisons, and no validated instruments, and the monograph will not pretend otherwise. What it takes from metrology is the \emph{grammar} of careful measurement talk, because the historical sources failed at exactly that grammar.

\subsection{Quantity, kind of quantity, and unit}\label{ch02:metrology:quantity}

In the Vocabulary's usage, a \emph{quantity} is a ``property of a phenomenon, body, or substance, where the property has a magnitude that can be expressed as a number and a reference'' \cite{vim}. Quantities partition into \emph{kinds}: length and wavelength are of the same kind; length and mass are not; and --- the point with teeth --- quantities of different kinds are not comparable, addable, or convertible merely because both are numbers. A \emph{unit} is a conventionally chosen reference quantity of the same kind, in terms of which magnitudes are expressed; a \emph{quantity value} is the pairing of a number with a unit.

Three consequences of this grammar recur throughout the monograph. First, a number without its unit and kind is not a quantity value at all; it is a numeral. Second, scales that happen to share a numeric range are not thereby the same kind: a score on 0--10 and a fraction on 0--1 differ in kind even where their numbers overlap, and a conversion between them is a substantive claim requiring a defined relation, not a rescaling one may perform silently. Third, ratios and products of quantities have derived kinds determined by their arguments; dividing a quantity of one kind by a quantity of another yields a third kind, which either has a place in the analyst's ontology or signals that the operation was meaningless. The Vocabulary also distinguishes \emph{ordinal} quantities --- those, like many scores, for which magnitude comparison is meaningful but differences and ratios are not --- and warns that arithmetic on ordinal quantities requires explicit justification. Score-valued quantities in any framework should be presumed ordinal until shown otherwise.

\subsection{Measurand, measurement, and calibration}\label{ch02:metrology:measurand}

The \emph{measurand} is the ``quantity intended to be measured'' --- and the Vocabulary is emphatic that specifying it is a substantive act: the ``same'' informal quantity under different conditions, definitions, or procedures is a different measurand \cite{vim}. A \emph{measurement} is the process of experimentally obtaining quantity values attributable to the measurand, and it presupposes a measurement procedure, a measuring system, and a description sufficient for another practitioner to repeat it. \emph{Calibration}, in the Vocabulary's sense, is the operation that establishes a relation between values indicated by a measuring system and values provided by standards, with their uncertainties, and then permits corrected measurement results; it is not a synonym for adjustment, tuning, or fitting in general.

The distinction between what a system \emph{indicates} and what one is licensed to \emph{report} as a measurement result is the deepest single lesson metrology offers a new field. An indication becomes a measurement result only through a chain: a defined measurand, a documented procedure, a calibration relation with uncertainty, and a stated coverage. Skipping the chain and reporting indications as results is the generic form of several failures the source audit catalogued.

\subsection{Uncertainty}\label{ch02:metrology:uncertainty}

The Guide's framework treats every measurement result as an estimate accompanied by a \emph{measurement uncertainty}: a ``non-negative parameter characterizing the dispersion of the quantity values being attributed to a measurand'' \cite{vim,gum}. Uncertainty evaluations are classified by method: \emph{Type~A} evaluations use statistical analysis of repeated observations (standard deviations of means, regression residuals); \emph{Type~B} evaluations use other knowledge --- manufacturer specifications, prior data, physical bounds, expert judgment --- encoded as probability distributions. Both types yield \emph{standard uncertainties}; components combine through the law of propagation of uncertainty (in first order, via sensitivity coefficients and covariances) into a \emph{combined standard uncertainty}, which may be scaled by a coverage factor into an \emph{expanded uncertainty} defining an interval with stated coverage probability.

Two disciplines from this framework are load-bearing later. First, \emph{uncertainty travels with the value}: a derived quantity's uncertainty is computed from its parents' uncertainties and covariances, not invented downstream, and a value whose uncertainty was never evaluated should be marked as such rather than silently treated as exact. Second, \emph{covariance is not optional}: when combined quantities share error sources, ignoring their covariance mis-states the combined uncertainty, sometimes drastically; where covariances are unavailable, the honest course is to declare the approximation being made rather than let a zero-covariance assumption pass as fact.

\subsection{Error, repeatability, and reproducibility}\label{ch02:metrology:error}

The Vocabulary draws a distinction, easily blurred in casual usage, between measurement \emph{error} --- the difference between a measured value and a reference value, a signed quantity that is in general unknowable when the true value is unknowable --- and measurement \emph{uncertainty} --- the analyst's quantified doubt, a non-negative dispersion parameter that is always evaluable \cite{vim,gum}. The older error-analysis tradition spoke of random and systematic errors; the uncertainty framework re-expresses the same physics without requiring knowledge of true values: systematic effects, once recognized, are corrected for, and the imperfect knowledge of the correction contributes a (Type~B) uncertainty component, while random effects contribute (Type~A) components estimated from dispersion. The practical rule is that recognized systematic effects are never left inside a reported value as silent bias; they are either corrected with an uncertainty on the correction or declared as a limitation.

Conditions of measurement are stratified with equal care. \emph{Repeatability} conditions hold everything fixed --- same procedure, same operators, same system, same conditions, same location, short time span --- and characterize the floor of dispersion attainable by the setup. \emph{Reproducibility} conditions vary what practice actually varies --- different laboratories, operators, instruments, occasions --- and characterize the dispersion a claim must survive to be portable. The gap between the two is itself diagnostic: results that repeat but do not reproduce indicate uncontrolled condition-dependence, and the metrological response is to promote the offending condition into the measurand's specification (a different condition then defines a different measurand) rather than to average over it silently. The transferable form of this rule appears throughout the candidate's design: configuration details that change results --- model revision, layer, direction provenance, corpus, seeds --- belong inside the profile identity, so that variation across them is variation across declared measurands, never noise smuggled into one.

\subsection{The calibration hierarchy}\label{ch02:metrology:hierarchy}

Institutional metrology arranges instrument calibrations in a \emph{hierarchy}: national standards realize units; reference standards are calibrated against them; working standards against those; and field instruments against working standards --- each step documented, each contributing uncertainty, so that the uncertainty of a field measurement provably includes every link above it \cite{vim}. Two structural features of the hierarchy are worth extracting from their institutional setting. First, uncertainty is \emph{monotone down the chain}: no link can report less uncertainty than its parent chain supplies, and any claimed shortcut (a field result more certain than its reference) is a red flag, not an achievement. Second, the hierarchy makes \emph{staleness} a first-class concept: calibrations expire, drift is monitored between calibrations, and a result traceable only through an expired or drifted link loses its status until recalibration. Both features transfer to any lineage-disciplined computational system as design requirements --- propagated uncertainty that can only grow through derivation, and expiry-and-drift metadata on every calibration-like binding --- and both appear, in exactly that role, in the candidate's calibration-profile contract discussed in later chapters.

\subsection{Traceability}\label{ch02:metrology:traceability}

\emph{Metrological traceability} is the ``property of a measurement result whereby the result can be related to a reference through a documented unbroken chain of calibrations, each contributing to the measurement uncertainty'' \cite{vim}. The definition is precise on three points: traceability is a property of \emph{results}, not of instruments or laboratories; the chain must be \emph{documented} and \emph{unbroken}; and every link must carry uncertainty. In institutional metrology the chain terminates in realizations of the SI units; the concept, however, is general --- a result is traceable to whatever reference the chain actually reaches, and no further.

The transferable content for a computational setting is the chain discipline itself: every reported value should be connected, by explicit and inspectable links, to the raw artifacts and procedures that produced it, with each link's method and uncertainty recorded. What does not transfer --- and must not be claimed --- is the institutional authority of the SI chain. A digest-linked lineage is documentation, not accreditation; a complete chain to a raw artifact establishes where a number came from, not that the number measures what its label says. That further step is validation, treated in Section~\ref{ch02:vv}.

\subsection{Missing values, censoring, and the discipline of absence}\label{ch02:metrology:censoring}

A final piece of measurement discipline concerns values that are \emph{not there}. Statistical practice distinguishes several kinds of absence, and the distinctions carry inferential weight. A value may be \emph{missing} --- the measurement was not made, or was made and lost --- in which case the salient question is whether the missingness is related to the value that would have been observed; analyses that ignore informative missingness are biased in ways no downstream care repairs. A value may be \emph{censored} --- the measurement was made but yields only partial information, as when a response falls below a detection limit (left-censoring) or an instrument clips at full scale (right-censoring); censored observations are not zeros, are not the limit value, and are not discardable: each of those common substitutions injects a specific, well-documented bias, and the correct treatments (likelihood contributions from the censoring interval, or explicit interval-valued records) are standard. And a value may be \emph{undefined} --- the quantity's defining ratio or operation has no value at the observed operating point, as when a ratio's denominator is zero; reporting a number anyway manufactures data.

The transferable discipline is representational: an evidence system must be able to say, of any field, ``absent, for this stated reason'' as a first-class recorded fact, distinct from zero, distinct from any sentinel numeral, and preserved through derivation --- so that a downstream consumer can distinguish a measured zero from a censored observation from an undefined quantity from a value never sought. This background point, elementary in biostatistics and survival analysis, is the rationale for the missing-or-censored-reason field of the candidate's observed-value record and for its censoring of residual diagnostics at zero-compatible response, both described in later chapters.

\section{Calibration of computer models}\label{ch02:calibration}

\subsection{The setting and the vocabulary collision}\label{ch02:calibration:setting}

A second, distinct sense of ``calibration'' governs the statistical literature on computer models. Here the object is a \emph{simulator}: a deterministic or stochastic code mapping inputs and unknown parameters to outputs, standing in for a physical or empirical system. \emph{Calibration} in this sense is the use of observed field data to learn about the unknown parameters --- and, crucially, about the simulator's inadequacy \cite{kennedyohagan}. The collision with the metrological sense of calibration (Section~\ref{ch02:metrology:measurand}) is unfortunate but entrenched; this monograph disambiguates by context and, where necessary, by the phrases \emph{instrument calibration} and \emph{model calibration}.

\subsection{Parameter uncertainty and model discrepancy}\label{ch02:calibration:kernel}

The canonical formulation writes the observed data as simulator output at the true-but-unknown parameter value, plus a \emph{model discrepancy} term, plus observation error:
\[
z_i = \rho\,\eta(x_i, \theta) + \delta(x_i) + \varepsilon_i,
\]
where $\eta$ is the simulator, $\theta$ the calibration parameters, $\delta(\cdot)$ an unknown function representing the systematic inadequacy of the simulator even at the best parameter value, and $\varepsilon_i$ observation error \cite{kennedyohagan}. Two commitments distinguish this framework from naive fitting. First, $\theta$ is treated as uncertain \emph{after} calibration: the analysis yields a posterior distribution, not a point estimate anointed as truth, and predictions propagate that posterior. Second --- the framework's signature --- the discrepancy term is carried explicitly: the model is assumed inadequate, the inadequacy is estimated alongside the parameters, and predictive uncertainty includes it. Omitting $\delta$ does not make a model adequate; it makes the analysis overconfident and biases the parameter estimates, which absorb the inadequacy in uninterpretable ways.

The framework also names an identifiability problem that any borrower must respect: with a flexible discrepancy, the data cannot always separate ``different parameter value'' from ``different inadequacy,'' so the physical interpretability of calibrated parameters is a claim requiring argument, not a free consequence of good fit. A calibrated parameter is, in the first instance, the parameter value that makes the simulator-plus-discrepancy predict well --- which is a weaker and more honest thing than the true value of a physical constant. The point has been sharpened since: carrying an explicit discrepancy term does not by itself restore interpretability, since with a weakly constrained $\delta$ the discrepancy absorbs error that would otherwise inform the parameters, and calibrated posteriors remain biased and overconfident; informative prior knowledge about the form of the inadequacy is required \cite{brynjarsdottir}.

\subsection{Train/held-out discipline and the circularity failure}\label{ch02:calibration:split}

Whatever the inferential machinery, the elementary discipline of predictive assessment is the separation of the data that estimate a model from the data that evaluate it. Assessing a fit on its own training data measures the model's capacity to absorb those data --- which increases with flexibility --- rather than its capacity to predict; the degenerate limit is the pure circularity in which a parameter is defined from a sample and then ``verified'' by reproducing that same sample exactly. Held-out evaluation, cross-validation, and prospective validation on later data are the standard remedies, each with known caveats (leakage between splits, repeated peeking, distribution shift between calibration and use).

For a proportionality model the degenerate limit is easily disguised, so it is written out here. Given one paired observation $(a, y)$ with $y \neq 0$, define $r := a/y$; then predict $\hat{y} := a/r$ and observe $\hat{y} = y$ exactly. The agreement is an algebraic identity, carries no information about any other dose, and survives every dimensional check. Its only reliable signature is procedural: the same sample appears on both sides. Hence the discipline --- registered splits, fit on one, evaluate on the other, report uncertainty from the fit --- is not methodological ornament but the minimal condition under which a fitted slope means anything at all. Later chapters encode this discipline structurally; the point here is that the discipline itself is standard, old, and non-negotiable.

\subsection{Emulation and computational practice}\label{ch02:calibration:emulation}

When the simulator is expensive, the calibration framework interposes an \emph{emulator}: a statistical surrogate --- in the canonical treatment, a Gaussian process --- trained on a designed set of simulator runs and used in place of the simulator inside the calibration computation, with the emulator's own approximation uncertainty propagated alongside parameter uncertainty and discrepancy \cite{kennedyohagan}. The details of emulation matter less for this monograph than the epistemic pattern it exemplifies: every layer interposed between data and conclusion --- surrogate, interpolation, preprocessing --- is itself a model with its own inadequacy, and a disciplined analysis carries that inadequacy explicitly rather than treating the layer as transparent. The pattern generalizes to computational measurement pipelines of any kind: a normalization step, a projection, a quantization each transform the quantity, and each transformation either has a declared method with an assessable error or constitutes an undeclared modeling assumption.

A second point of computational practice: \emph{the design of calibration runs is part of the inference}. Which parameter settings are simulated, over what ranges, with what space-filling properties, determines what the calibration can learn --- exactly parallel to the dose-grid point of Section~\ref{ch02:dose:design}. Sparse or lopsided designs yield posteriors that reflect the design as much as the system, and the standard practice is to report the design alongside the result.

\subsection{Validation of calibrated models and extrapolation}\label{ch02:calibration:validation}

Calibration and validation interact in a way the literature is careful about: data used to calibrate cannot also validate, since the calibrated model has been conformed to them. Validation of a calibrated model therefore requires either held-out field data, prospective data gathered after freezing the model, or structured cross-validation with the calibration repeated inside each fold \cite{sargent}. Moreover, a calibrated-and-validated model is validated \emph{over the domain of the validation data}; the discrepancy term, being an empirical estimate of inadequacy where data exist, provides no warrant where they do not, and extrapolation beyond the calibration-and-validation envelope reintroduces unquantified inadequacy. The framework thus arrives, by its own route, at the same regime-boundedness that causal inference imposes (Section~\ref{ch02:causal}): a fitted, uncertainty-carrying, discrepancy-aware model earns claims inside its envelope and speculation outside it. Later chapters implement this as an expiry-and-envelope discipline on calibration profiles; the principle recorded here is standard.

\section{Verification and validation}\label{ch02:vv}

\subsection{Two questions that must not be merged}\label{ch02:vv:questions}

The simulation literature distinguishes two questions about any model-bearing artifact. \emph{Verification} asks: was the artifact built correctly --- does the implementation faithfully realize its specification, free of the defects that testing can expose? \emph{Validation} asks: was the correct artifact built --- does the model represent the target system accurately enough for its intended use \cite{sargent}? The questions have different evidence types, different audiences, and different failure modes, and a mature practice keeps them separate: a model can be perfectly verified and wholly invalid (a flawless implementation of a wrong idea), or roughly right in conception yet unusable because its implementation is untrustworthy.

Verification evidence is the familiar apparatus of software quality: unit and property tests, boundary-case fixtures, golden rejection tests (inputs that \emph{must} fail, verifying that guards guard), replay of pinned computations, static analysis. Validation evidence is comparison against the world: agreement with independent data not used in construction, behavior under conditions spanning the intended operating envelope, expert scrutiny of assumptions, and --- in strong practices --- validation performed or reviewed by parties independent of the developers \cite{sargent}. The literature further ties validity to \emph{intended use}: a model is not valid simpliciter but valid for a stated purpose over a stated domain, and reuse outside that domain restarts the validation question.

\subsection{A catalogue of validation techniques}\label{ch02:vv:techniques}

The simulation literature has accumulated a standard catalogue of validation techniques, and rehearsing a representative subset clarifies how varied --- and how demanding --- validation evidence is compared with verification evidence \cite{sargent}. \emph{Comparison to other models}: outputs are checked against independently developed models of the same system, valued precisely because shared errors between independent efforts are less likely. \emph{Degenerate tests}: parameters are pushed to values with known degenerate behavior (zero arrival rate, infinite capacity) and the model is required to reproduce the known limit. \emph{Extreme-condition tests}: the model must remain plausible at the boundaries of its input space, where errors concentrate. \emph{Event and face validity}: occurrence patterns and input-output behavior are reviewed by people knowledgeable about the system for gross implausibility --- weak evidence alone, but cheap and early. \emph{Historical data validation}: part of the historical record drives the model, the remainder tests it --- the split discipline of Section~\ref{ch02:calibration:split} in its simulation form. \emph{Parameter variability and sensitivity analysis}: inputs are varied and the direction and magnitude of output changes checked against system knowledge; parameters found sensitive demand accurate estimation. \emph{Predictive validation}: the model forecasts, and the system is then observed. \emph{Internal validity}: replicated stochastic runs are checked for a stability of results without which no single run means anything. \emph{Traces}: individual entities are followed through the model logic to confirm the implemented behavior is the intended one --- a technique straddling the verification-validation border.

Two features of the catalogue merit emphasis. First, no single technique validates; practices accumulate evidence across techniques, matched to the model's intended use and the cost of error. Second, several techniques presuppose access to the target system --- historical data, predictive trials, expert knowledge of system behavior --- and are simply unavailable when no comparable system data exist; a practice honest about this reports which techniques were unavailable rather than substituting the available ones and claiming completeness.

\subsection{Documentation, independence, and accreditation}\label{ch02:vv:independence}

The literature further separates the \emph{performance} of verification and validation from its \emph{evaluation}. Documentation --- of assumptions, techniques applied, results, and known limitations --- is treated as constitutive: undocumented validation is unevaluable and therefore, for consumers of the model, indistinguishable from no validation. Independence enters in degrees: developer self-assessment is the weakest arrangement; independent verification and validation by parties without a stake in the model's acceptance is materially stronger; and formal \emph{accreditation} --- an authority's determination that a model is acceptable for a specific purpose --- is an institutional act distinct from both, performed against declared criteria \cite{sargent}. The mapping onto this monograph's status vocabulary is direct and is stated once here: everything the candidate's own project produces, including this monograph, is developer self-assessment; the program's internal cross-review would constitute a first independence step; external independent review a second; and nothing corresponding to accreditation exists or is claimed. The status line \textsc{open / ring 4 / not licensed / held} is, among other things, a compact encoding of where in this hierarchy the candidate currently sits: at the bottom.

\subsection{Operational consequence for this monograph}\label{ch02:vv:consequence}

The consequence for reading everything that follows is mechanical. Test-suite results --- however extensive --- are verification: they establish that specified behavior, including specified rejection behavior, occurs at a pinned revision on constructed inputs. They cannot establish that the specification captures anything true about steered neural activations; that is validation, requires the falsification program's controlled runs against held-out data, and has not occurred. When later chapters cite passing tests over synthetic known-fault fixtures, the citation licenses exactly the sentence ``the software detects the synthetic fault it was built to detect'' and no sentence about live models. This is the standard reading \cite{sargent}, stated once here so that later citations can be brief.

\section{Limits of causal inference from observational fits}\label{ch02:causal}

\subsection{The ladder}\label{ch02:causal:ladder}

Modern causal inference, in the structural tradition, organizes claims into three strictly ordered levels: \emph{association} (what the joint distribution of observed variables looks like), \emph{intervention} (what would happen to one variable if another were set by external action), and \emph{counterfactuals} (what would have happened in a particular case had circumstances differed) \cite{pearl}. The central formal lesson is that these levels do not reduce to one another: no amount of associational data, by itself, answers interventional questions, because distinct causal structures can generate identical joint distributions. Interventional content enters only through interventional assumptions (encoded, e.g., as a structural causal model or graph) or interventional data (actual experiments), and counterfactual content requires still more.

The apparatus that makes this precise includes the do-operator, distinguishing $P(y \mid \mathrm{do}(a))$ --- the distribution of $y$ under external setting of $a$ --- from the observational $P(y \mid a)$; graphical criteria for when interventional quantities are \emph{identifiable} from observational ones given an assumed graph; and the systematic study of \emph{confounding}, the presence of common causes that induce association without causation \cite{pearl}.

\subsection{What a fitted slope is and is not}\label{ch02:causal:slope}

Where the dose is genuinely randomized and externally set by the experimenter, a dose-response fit does estimate an interventional quantity --- for the manipulated variable, in the sampled population, over the sampled range. This is the classical warrant of the randomized experiment, and it is why the falsification program's registered designs randomize the alpha dose and include controls. But the warrant is exactly as wide as the randomization and no wider, and three restrictions matter even in the best case.

First, \emph{regime-boundedness}: an interventional relation estimated over a grid supports claims about doses within (approximately) that grid, under that protocol, on that population; extrapolation beyond the grid or transport to a different population are additional assumptions, not consequences \cite{pearl}. Second, \emph{mechanism-silence}: even a well-estimated interventional slope says nothing about \emph{how} the effect is mediated; competing mechanistic stories --- direct paths, mediated paths, multiple parallel paths --- can share the same low-dimensional interventional summary, and distinguishing them requires interventions on the mediators themselves. Third, and most relevant to the candidate's lineage, \emph{residual non-identification}: the deviation of an observation from a fitted prediction is a composite of model misspecification, measurement error, unmodeled heterogeneity, and --- possibly --- unmodeled causes. Attributing a residual's sign to a specific hidden cause is a causal claim at the interventional or counterfactual level, and it is not identified by the fit; making it honestly requires an explicit causal model in which the hidden cause appears, plus interventions (e.g., deliberately implanted faults with controls) that could distinguish it from the alternatives \cite{pearl}. The audit's residual conflict, and the candidate's censoring discipline, are downstream of exactly this point.

\subsection{Adjustment and its perils}\label{ch02:causal:adjustment}

Between the poles of pure observation and randomized intervention lies the everyday practice of \emph{adjustment}: regressing an outcome on a treatment while controlling for covariates, in the hope that conditioning removes confounding. The structural literature gives this practice both its warrant and its boundary \cite{pearl}. The warrant is the identification result: if a covariate set satisfies the appropriate graphical criterion relative to an assumed causal graph --- blocking the spurious (back-door) paths between treatment and outcome without opening others --- then the adjusted association equals the interventional quantity, and observational data suffice. The boundary is everything the warrant presupposes: the graph is an assumption, not an observation; omitting a genuine confounder leaves bias in place; and, less intuitively, conditioning can \emph{create} bias --- adjusting for a variable that is a common effect of treatment and outcome (a collider), or a consequence of treatment on the causal path, induces spurious association or blocks part of the effect being estimated. ``Control for everything measurable'' is therefore not a conservative strategy but an incoherent one; which covariates to adjust for is determined by causal structure, and without a defended structure, adjusted estimates have no defended interpretation.

The lesson is recorded here because steering-style analyses invite exactly this practice: comparing fitted parameters across layers, directions, prompts, or checkpoints while ``controlling'' for whatever covariates are convenient. The background licenses such comparisons as descriptive summaries only; elevating any of them to a causal contrast would require the graph-level argument this section describes, which no part of the candidate supplies or claims.

\subsection{Observational overlays}\label{ch02:causal:overlays}

Finally, not every quantity in a steering-style experiment is randomized even when the dose is: the choice of layer, direction, corpus, and test point is typically made adaptively, on the basis of prior observation, and comparisons \emph{across} those choices are observational comparisons subject to selection and confounding in full force. A fitted parameter that predicts well within its profile carries no license to explain \emph{why} one profile differs from another. The general rule --- interventional claims extend exactly as far as the interventions --- covers both this and the preceding restrictions, and it is the rule later chapters encode.

\section{Activation steering and representation engineering}\label{ch02:steering}

\subsection{Status of this section}\label{ch02:steering:status}

This section surveys, at primer depth, the empirical literature within which a CAS-Ohm experiment would be situated: activation steering and, more broadly, representation engineering in neural networks. The status rules of Section~\ref{ch02:purpose} apply with maximum force here. What follows describes what published work in this area does and the kinds of phenomena it reports \cite{zourepe}; it does not assert that those phenomena are settled, that they transfer to any particular model or configuration, or --- above all --- that any CAS-Ohm measurement has been performed or any CAS-Ohm hypothesis supported. The candidate treats every quantitative question about steering in its own setting as open.

\subsection{Activations, directions, and readouts}\label{ch02:steering:setting}

A transformer-style network computes, at each layer and token position, a high-dimensional activation vector; the sequence of these vectors along the depth of the network (often described via the residual stream, the running sum to which each block adds its contribution) constitutes the network's internal state for a given input. Interpretability research has long explored the hypothesis that behaviorally meaningful features are represented, at least in part, as \emph{directions} in these activation spaces: that movement of the activation along a particular vector corresponds to a change in some property of the computation or the output.

Two operational primitives follow from this picture and recur throughout the literature. The first is the \emph{directional readout}: given a unit-norm direction $v$ at a declared layer and position, the scalar projection $\langle h, v \rangle$ of the activation $h$ onto $v$ is recorded as a measurement of ``how much'' of the putative feature is present. The second is the \emph{activation intervention}: the activation is modified in place during the forward pass --- most simply by adding a scaled direction, $h \mapsto h + \alpha v$ --- and the effect on downstream computation or generated text is observed. The scalar $\alpha$ is the steering coefficient: its sign selects the direction's orientation, and its magnitude the strength of the intervention.

\subsection{Steering vectors}\label{ch02:steering:vectors}

The activation-addition literature constructs steering vectors in several ways, of which four families suffice for a primer. \emph{Difference-based} constructions take the difference of activations induced by contrasting inputs --- for instance, a pair of prompts embodying opposite values of the property of interest --- optionally averaged over many pairs, on the reasoning that the difference cancels shared content and retains the contrast \cite{zourepe}. \emph{Probe-based} constructions train a linear classifier (a probe) to predict the property from activations and take the probe's weight vector as the direction. \emph{Optimization-based} constructions search directly for a vector whose addition produces a target behavior. \emph{Decomposition-based} constructions take as the direction a feature learned by an unsupervised decomposition of the activation space --- most prominently a sparse-autoencoder feature --- and steer along it \cite{templeton}; benchmark comparison reports that such feature steering, like the other families, can underperform even simple prompting baselines on behavioral steering tasks \cite{axbench}. Each family carries characteristic caveats discussed in the literature: difference vectors inherit whatever else differs between the contrasting inputs; probe directions reflect what is \emph{decodable}, which need not be what is \emph{used} by the network; optimized vectors can exploit idiosyncrasies of the objective. Across families, reported practice varies in normalization (unit norm or not), in the layer and token positions targeted, and in whether the same vector is applied at one position or broadcast across many --- all of which are consequential experimental parameters rather than nuisance details.

\subsection{Interventions at layers, and coefficient sweeps}\label{ch02:steering:sweeps}

The intervention's locus --- which layer, which positions, which hook point within the block --- is empirically consequential in the reported literature: the same nominal direction applied at different depths can produce effects of different size or character, and part of experimental practice is sweeping or selecting the layer \cite{zourepe}. Equally standard is the \emph{coefficient sweep}: applying the intervention over a grid of $\alpha$ values and recording, at each, either a downstream behavioral measure or an internal readout. Coefficient sweeps in this tradition can exhibit a range of qualitative shapes --- regions where effects grow with coefficient, regions where output degrades into disfluency at large magnitude, asymmetries between positive and negative coefficients --- and the literature is candid that evaluation is difficult: behavioral effects are measured by classifiers, human raters, or model-based judges, each with their own error, and internal readouts measure the projection, not the behavior. The primer-level summary is that the coefficient sweep is exactly a dose-design (Section~\ref{ch02:dose:design}) applied to activation space, with all of that tradition's questions --- where is the informative range, what does flatness mean, when does the preparation break --- reappearing in new clothes. Whether any given configuration exhibits a clean region anatomy is, in each case, an empirical finding of that experiment, not a property of the method.

\subsection{Controls in the steering setting}\label{ch02:steering:controls}

Because a steering intervention changes the activation by construction, the inferentially hard question is never whether \emph{something} changed but whether the change is attributable to the direction's semantic content rather than to generic perturbation. The literature's response is a family of controls, which this survey records as standard practice: \emph{zero-intervention} runs establishing the baseline; \emph{random-direction} controls of matched norm, testing whether an arbitrary vector of the same size produces comparable effects; \emph{wrong-layer} and \emph{wrong-position} applications, testing locality; \emph{shuffled or mismatched} pairings between inputs and vectors; and norm-matching generally, since effect size scales with perturbation magnitude and unmatched comparisons confound content with size. The methodological point --- a difference measured against no control, or against an unmatched control, does not support a content-specific claim --- is standard experimental logic; the steering literature's contribution is adapting that logic to activation space \cite{zourepe}.

\subsection{Evaluation difficulties and reproducibility cautions}\label{ch02:steering:evaluation}

The steering literature is, at its best, candid about the difficulty of evaluating its own interventions, and a primer should transmit the cautions along with the techniques \cite{zourepe}. The central difficulty is that the quantity of ultimate interest --- a change in model \emph{behavior} along a semantic dimension --- is not directly readable from any scalar. Behavioral evaluation routes through instruments that are themselves models or judgments: trained classifiers over outputs (with their own error rates and blind spots), human raters (with inter-rater disagreement and prompt-order effects), or judge models (whose failure modes may correlate with the very interventions under study). Internal evaluation routes through readouts like directional projections, which are cheap and precise but measure the representation, not the behavior; the coupling between a projection shift and a behavioral shift is itself an empirical hypothesis, not an identity. A result established under one evaluation instrument does not automatically hold under another, and the literature's stronger studies report multiple instruments and their disagreements.

Reproducibility cautions compound the evaluation difficulty. Reported steering effects are, in the published record, sensitive to experimental parameters that casual summaries omit: the model checkpoint (effects need not survive fine-tuning or even re-runs of training), the tokenizer's segmentation of the probe prompts, the layer and position conventions (off-by-one conventions differ across codebases), the norm convention for the direction, and the sampling temperature of generation-based evaluations. Sign conventions are a particular hazard: a direction and its negation are equally valid outputs of a difference or probe construction, and cross-study comparison of ``positive'' steering presupposes a shared orientation convention that must be checked, not assumed. None of these cautions is a criticism of the literature; they have since been documented systematically: large-scale evaluations report that steerability is highly variable across inputs and brittle to reasonable prompt changes \cite{tansteering}, and that prominent steering methods fail to improve --- and at times degrade --- behaviour on a substantial fraction of models tested \cite{dasilva}. The cautions are recorded here because they translate directly into requirements on any measurement ontology that would operate in this setting --- checkpoint digests, tokenizer digests, explicit hook and position identities, declared norms and signs --- requirements whose implementation in the candidate is described in Chapter~\ref{ch03:ontology}.

\subsection{Representation engineering more broadly}\label{ch02:steering:repe}

Activation addition is one member of a broader family sometimes collected under the name representation engineering: techniques that read or write network-internal representations to characterize or influence behavior. On the reading side, the family includes linear probing at scale, concept-activation analyses, and unsupervised decompositions of activation spaces; on the writing side, ablations that remove a direction (projecting activations onto its orthogonal complement), targeted patches that transplant activations between runs, and learned interventions trained against behavioral objectives. The family matters to this survey for one structural reason: it supplies \emph{alternative operationalizations} of what is informally the same manipulandum. A ``feature'' can be operationalized as a probe direction, a difference direction, a decomposition component, or an ablation target, and these operationalizations need not agree --- a probe may decode what an ablation shows to be causally inert, and a steering vector may act through content no probe was trained to see. The metrological translation (Section~\ref{ch02:metrology:measurand}) is exact: each operationalization defines a different measurand, and results attach to the measurand actually used, not to the informal feature name. The candidate's insistence that direction \emph{provenance} --- probe-derived, difference-derived, or otherwise --- be part of every test-point identity is the mechanical form of this point.

\subsection{What this setting supplies and what it leaves open}\label{ch02:steering:open}

For this monograph, the steering literature supplies the empirical coordinates a response-model candidate needs: a controllable scalar dose (the steering coefficient), a measurable scalar response (a directional projection at a declared test point), an inventory of consequential experimental parameters (model identity, layer, positions, hook, direction provenance, norm, sign, corpus), and a control repertoire. It leaves open --- as questions untouched in this project, and only partially answered in that literature, where systematic studies report that steering effects transfer unreliably across prompts and settings \cite{tansteering,dasilva} --- whether projection responses to coefficient sweeps are lawful in any reusable sense, whether fitted summaries transfer across prompts, layers, or checkpoints, and whether any composition rules hold. The candidate's registered hypotheses are bets on some of these questions; this chapter, being background, places no bets.

\section{Provenance}\label{ch02:prov}

\subsection{The PROV model}\label{ch02:prov:model}

The W3C provenance recommendations define a domain-neutral model for recording how things came to be. Its core ontology, PROV-O, is built from three node classes and a small set of relations \cite{prov}. \emph{Entities} are things --- documents, datasets, records, values. \emph{Activities} are occurrents that use and generate entities over time. \emph{Agents} are things bearing responsibility --- persons, organizations, software. The load-bearing relations include: an entity \emph{was generated by} an activity; an activity \emph{used} entities; an entity \emph{was derived from} another entity; an entity or activity \emph{was attributed to} or \emph{associated with} an agent; and an agent \emph{acted on behalf of} another. From these primitives, arbitrarily deep lineage graphs can be assembled: every value can point to the activities that produced it, the entities those consumed, and the agents responsible, recursively to the raw sources.

Provenance modeling of this kind is what makes lineage \emph{inspectable}: given a questioned result, one can walk the graph to its inputs, identify every transformation, and ask of each link whether it was appropriate. It underwrites reproducibility (the graph documents what to redo), attribution (the graph names who and what is responsible), and impact analysis (the graph shows what downstream results a defective input contaminates).

\subsection{Qualified relations, bundles, and plans}\label{ch02:prov:qualified}

Beyond its core triangle, the PROV model supplies refinements that recur in serious lineage systems \cite{prov}. \emph{Qualified relations} expand a bare edge into a described one: instead of merely asserting that an entity was derived from another, a qualified derivation can name the activity that effected it, the specific usage and generation events involved, and the time; instead of merely associating an activity with an agent, a qualified association can name the \emph{plan} --- the procedure or protocol --- the agent was following. \emph{Roles} type the manner of participation: the same activity may use one entity as input data and another as configuration, and the role distinction preserves which was which. \emph{Bundles} make provenance itself first-class: a bundle is a named set of provenance assertions whose own provenance --- who asserted this history, when, following what --- can then be described, enabling meta-lineage and the auditing of audits. Collections and their memberships allow lineage over aggregates without losing the identities of members.

These refinements matter here because the naive reading of provenance --- a bare ancestry graph --- is too weak for measurement discipline. A measurement's lineage must record not just \emph{that} a value descends from raw artifacts but \emph{by what method}, \emph{under which configuration}, \emph{in what role} each parent participated, and \emph{in which order} --- distinctions that map onto qualified derivations, plans, and roles. The candidate's typed parent references (role-tagged, position-ordered, digest-identified) and method identities are, in PROV terms, qualified derivations under declared plans; the correspondence is noted in Chapter~\ref{ch03:ontology}, and recorded here as an application of the standard model rather than a new idea.

\subsection{Structure is not truth}\label{ch02:prov:truth}

The recommendation itself is careful about what provenance establishes, and the caution is restated independently here because it bounds a hazard this monograph has already flagged (Chapter~\ref{ch01:intro}, authority boundary). A provenance record is an \emph{assertion} about history, made by some agent; its existence establishes that the assertion was made, not that it is true, and a well-formed graph can describe a fabricated history as fluently as a real one \cite{prov}. Trust in provenance therefore requires machinery beyond structure: integrity protection (so records cannot be silently altered), authentication (so asserting agents are who they claim), and review (so the asserted history is checked against independent evidence). Cryptographic digests and signatures supply the first and part of the second; nothing in the provenance layer supplies the third. The transferable rule --- \emph{provenance structure is documentation, never authority} --- is the background form of the candidate's split between structural completeness and authenticated eligibility, and it is why that split cannot be closed by accumulating more structure.

\section{Synthesis: what the background licenses}\label{ch02:synthesis}

Assembled, the seven literatures surveyed here license a specific, limited methodological package, which the remainder of the monograph draws on by reference.

From dose-response analysis: the dose/response framing, the region anatomy as signature classification, locality of linear summaries, and the flat-response ambiguity \cite{ritz}. From metrology: the quantity/kind/unit grammar, the measurand's procedure-dependence, uncertainty as an obligatory companion of every value with propagation and declared covariance treatment, and chain-style traceability as documentation \cite{vim,gum}. From model calibration: parameter uncertainty after fitting, explicit model discrepancy, and the identifiability caution on interpreting calibrated parameters \cite{kennedyohagan,brynjarsdottir}, together with the train/held-out discipline whose violation is circularity. From verification and validation: the strict separation of built-correctly from built-the-correct-thing, and the confinement of test-suite evidence to the former \cite{sargent}. From causal inference: the association/intervention/counterfactual hierarchy, regime-bounded warrants for randomized doses, mechanism-silence of interventional summaries, and the non-identification of causes from residuals \cite{pearl}. From the steering literature: the empirical coordinates --- coefficients, directions, projections, layers --- their consequential experimental parameters, and the control repertoire \cite{zourepe}. From provenance: the entity/activity/agent lineage model and the rule that structure is not truth \cite{prov}.

Table~\ref{ch02:tab:map} summarizes the mapping from each background discipline to the concepts this monograph draws from it and the locus of their later application. The table is a navigation aid, not an argument; each row's content is exactly what the corresponding section stated, no more.

\begin{table}[ht]
\centering
\small
\begin{tabular}{@{}lll@{}}
\toprule
Discipline & Concepts drawn & Later locus \\
\midrule
Dose-response \cite{ritz} & regions; local linearity; flat ambiguity & region axis; active-region gates \\
Metrology \cite{vim,gum} & kind/unit grammar; uncertainty; & quantity kinds; unit ontology; \\
 & traceability; censoring discipline & lineage; censored fields \\
Model calibration \cite{kennedyohagan} & parameter uncertainty; discrepancy; & profile contract; fit roles; \\
 & split discipline & anti-circularity rule \\
V\&V \cite{sargent} & verification/validation split; & bounded software evidence; \\
 & independence hierarchy & status vocabulary \\
Causal inference \cite{pearl} & intervention hierarchy; adjustment & residual censoring; claim \\
 & limits; residual non-identification & boundaries \\
Steering \cite{zourepe} & doses, directions, projections; & test-point identity; controls \\
 & controls; parameter sensitivity & in registered hypotheses \\
Provenance \cite{prov} & lineage model; qualified relations; & parent refs; method identities; \\
 & structure is not truth & eligibility split \\
\bottomrule
\end{tabular}
\caption{Background disciplines, the concepts this monograph draws from each, and where the candidate applies them. Every row is background applied, not novelty claimed.}
\label{ch02:tab:map}
\end{table}

What the background does not license is equally definite. It does not license any claim that steered activations obey a proportional law --- the dose-response frame classifies signatures; it does not guarantee their existence. It does not license treating scores as measurements, indications as results, or lineage as authority. It does not license causal readings of fitted slopes beyond their randomized, regime-bounded warrant, nor any causal reading of residuals. It does not license transporting steering phenomena reported elsewhere into claims about the candidate's own setting. Every empirical sentence about CAS-Ohm 1.0 therefore remains exactly where Chapter~\ref{ch01:intro} left it: registered, controlled for, and open. The next chapter turns from the literatures to the candidate's first contribution --- the typed ontology that makes this discipline mechanical.
